%% ============================================================
%% Einleitung
%% ============================================================

\chapter{Einleitung}
\label{chap:einleitung}

% Hinweis: Die Einleitung fuehrt in das Thema ein und umfasst
% typischerweise 2-3 Seiten. Sie enthaelt Hinfuehrung zum Thema,
% Forschungsfrage, Methodik und Aufbau der Arbeit.

[Hinfuehrung zum Thema: Warum ist das Thema relevant? Welche aktuelle
Bedeutung hat es? Hier wird das Interesse der Leserin / des Lesers geweckt
und der Kontext der Arbeit hergestellt.]

[Persoenliche Motivation oder gesellschaftliche Relevanz des Themas
erlaeutern.]


\section{Forschungsfrage}
\label{sec:forschungsfrage}

Die zentrale Forschungsfrage dieser Vorwissenschaftlichen Arbeit lautet:

\begin{quote}
\textit{[Hier die Forschungsfrage als vollstaendigen Fragesatz formulieren,
\zB \glqq Inwiefern beeinflusste die Industrialisierung die soziale Struktur
in Wiener Neustadt im 19.~Jahrhundert?\grqq]}
\end{quote}

% Optional: Untergeordnete Fragestellungen
Zur Beantwortung dieser Fragestellung werden folgende Teilfragen untersucht:
\begin{itemize}
  \item [Teilfrage 1]
  \item [Teilfrage 2]
  \item [Teilfrage 3]
\end{itemize}


\section{Methodik}
\label{sec:methodik}

Zur Beantwortung der Forschungsfrage wird in dieser Arbeit
[Methodik beschreiben] eingesetzt. [Erklaeren, warum diese Methode
gewaehlt wurde und wie sie angewendet wird.]

% Beispiele fuer Methoden:
% - Literaturanalyse: Auswertung von Fachliteratur und Primaerquellen
% - Qualitative Interviews: Leitfadengestuetzte Interviews mit ExpertInnen
% - Vergleichende Analyse: Gegenueberstellung von Fallbeispielen
% - Empirische Erhebung: Fragebogen, Beobachtung, Experiment

[Konkrete Beschreibung der Vorgehensweise, verwendete Quellen und
Materialien.]


\section{Aufbau der Arbeit}
\label{sec:aufbau}

Die vorliegende Arbeit gliedert sich in [Anzahl] Kapitel.

Nach dieser Einleitung werden in Kapitel~\ref{chap:theoretische-grundlagen}
die theoretischen Grundlagen dargestellt. Kapitel~\ref{chap:analyse}
widmet sich [Beschreibung]. In Kapitel~\ref{chap:ergebnisse} werden
die Ergebnisse praesentiert. Abschliessend fasst Kapitel~\ref{chap:fazit}
die zentralen Erkenntnisse zusammen und beantwortet die Forschungsfrage.
